\documentclass{tpks}
\doctitle{Системные требования}
\docversion{1}
\docdate{03.10.2026}

\begin{document}
\makeheader

\section{Введение}
\subsection{Постановка задачи}

Релиз программного продукта объединяет версии нескольких компонентов. Перед выпуском
эти версии проходят установку и проверки в нескольких средах. Для решения о выпуске
нужны сведения о составе релиза, результатах проверок, согласованиях и текущих установках.

В проекте разрабатывается система автоматизации управления жизненным циклом релиза ПО
«Диспетчер релизов». Она хранит эти сведения и управляет прохождением этапов.
Сборку, установку и автоматические проверки выполняют внешние системы.

\subsection{Назначение и цели системы}

Система предназначена для подготовки релиза и его последовательного продвижения до
промышленной среды. Цели: сохранить известный состав релиза на всём маршруте,
допускать его к этапам по установленным правилам, учитывать установки и предоставлять
историю действий. При необходимости система организует откат к предыдущему выпуску.

Документ задаёт поведение системы в целом. На его основе будут разработаны требования
к ПО и системные испытания. Формулировки с ID в разделах~\ref{sec:functional}
и~\ref{sec:nonfunctional} являются требованиями. Сценарии поясняют их применение.

\subsection{Область действия}

Версия 1 рассчитана на одну организацию, несколько программных продуктов и семь
фиксированных ролей пользователей. Один пользователь может иметь несколько ролей.
Для каждого продукта задаётся один маршрут; его этапы и шаги выполняются последовательно.
Предлагаемая нагрузка --- 20 одновременно продвигаемых релизов и 30 одновременно
работающих пользователей. Эти значения требуют подтверждения по пункту В-5
раздела~\ref{sec:questions}.

В границы системы входят каталог, состав релизов, маршруты, допуски к этапам,
учёт установок, организация отката, права и аудит. Исходный код и дистрибутивы хранятся
внешними средствами. Роли пользователей относятся к работе продукта;
состав учебной команды указан в шапке документа.

\subsection{Термины и сокращения}

Определения взяты из общего глоссария проекта. Роли определены отдельно в
разделе~\ref{sec:roles}. Обозначение П-N используется для потребности,
ВИ-N --- для варианта использования, В-N --- для открытого вопроса.
ID означает идентификатор; SYS обозначает системное требование.
Единицы измерения: с --- секунда, \% --- процент.

\printterms{Аутентификация, Версия компонента, Внешняя система развёртывания,
Внешняя система сборки, Внешняя система учёта задач, Внешняя служба аутентификации,
Внешняя служба доставки сообщений, Выпуск релиза, Дистрибутив, Журнал аудита,
Компонент, Контрольная сумма, Критический дефект, КТК, Маршрут продвижения,
Откат релиза, Перечень изменений релиза, ПО, Правило согласования,
Принцип «четырёх глаз», Программный продукт, Продвижение релиза, Промышленная среда,
ПСИ, Ревизия исходного кода, Релиз, Релиз-кандидат, Сборка, Согласование,
Состав релиза, Среда развёртывания, Технологическое окно, Установка,
Фиксация состава релиза, Хранилище дистрибутивов, Шаг, Этап}

\subsection{Ссылки}

\begin{enumerate}
\item Козырев В.~П. Слайды «Требования» и «Управление конфигурацией ПО»:
\path{materials/kozyrev_slides/Vse_kozyrev.pdf}, с.~16--27, 31, 78.
Использованы правила описания, проверки и идентификации требований.
\item \href{../materials/lectures/lectures_konspekt.txt}{Конспект лекций Б21-504},
строки 38--67, 106, 161.
Использован как дополнительный источник; при расхождениях приоритет имеют слайды.
\item Рабочие материалы команды: \path{notes/04_dpm_platform_v.md}, часть Б,
разделы 1--10. Источник потребностей, модели предметной области и черновика требований.
\item Platform V DevOps Pipeline Management (DPM), документация версии 4.10.0:
\href{https://platformv.sbertech.ru/docs/public/DPM/4.10.0/common/index.html}{официальная документация}.
Референс предметной области; требования к учебной системе сформулированы командой.
\end{enumerate}

\section{Общее описание}
\subsection{Окружение системы}

Пользователи задают состав релиза, маршрут и условия допуска, запускают продвижение,
принимают решения и просматривают результаты. Внешняя система сборки передаёт сведения
о версиях. Система развёртывания получает задания и возвращает результаты.
Система учёта задач предоставляет перечень изменений и сведения о дефектах.
Служба аутентификации подтверждает личность пользователя.

Пятая внешняя система --- служба доставки сообщений --- включена в описание окружения
как резервное расширение. Обязательные требования версии 1 не зависят от доставки
уведомлений; назначенные действия доступны пользователям в самой системе.
Состав обмена приведён в таблице~\ref{tab:interfaces}.

\subsection{Роли пользователей}\label{sec:roles}

Права перечислены в таблице~\ref{tab:roles}. Всем семи ролям доступен просмотр
каталога, релизов, их состояния и установок. Для остальных операций действует
указанное распределение. Совмещение ролей даёт объединение прав, но сохраняет
запрет согласовывать собственное продвижение.

\begin{tabl}{|P{3.5cm}|L|}{Роли и разрешённые действия}[tab:roles]{\textbf{Роль} & \textbf{Действия}}
Администратор & Регистрирует продукты, компоненты и среды; назначает роли пользователям. \\ \hline
Релиз-менеджер & Создаёт релиз, меняет и фиксирует его состав; настраивает маршрут, КТК и согласования; запускает продвижение, повторяет шаг, отменяет релиз и инициирует откат. \\ \hline
Разработчик & Просматривает каталог, релизы, результаты продвижения и установки. \\ \hline
Инженер по тестированию & Регистрирует результат назначенного ему ручного шага на непромышленном этапе. \\ \hline
Инженер эксплуатации & Задаёт технологические окна; регистрирует результат назначенного ему ручного шага в промышленной среде. \\ \hline
Согласующий & Принимает решение по адресованному ему запросу на согласование. \\ \hline
Аудитор & Просматривает журнал аудита. \\ \hline
\end{tabl}

\subsection{Потребности}

Потребности таблицы~\ref{tab:needs} служат источниками требований.

\begin{tabl}{|l|L|}{Потребности пользователей}[tab:needs]{\textbf{ID} & \textbf{Потребность}}
П-1 & Знать состав каждого релиза и включённые в него изменения. \\ \hline
П-2 & Проводить релиз через среды в установленном порядке. \\ \hline
П-3 & Допускать релиз к следующему этапу после проверки качества и согласования. \\ \hline
П-4 & Исключать конфликт работ в общей среде и учитывать разрешённое время начала работ. \\ \hline
П-5 & Возвращать промышленную среду к предыдущему выпуску при проблемах с новым. \\ \hline
П-6 & Знать, какие версии компонентов установлены в каждой среде. \\ \hline
П-7 & Иметь историю изменений, защищённую от правки задним числом. \\ \hline
П-8 & Выполнять действия только с правами соответствующей роли. \\ \hline
П-9 & Использовать сведения и результаты внешних систем сборки, развёртывания и учёта задач. \\ \hline
\end{tabl}

\subsection{Варианты использования}

\textbf{ВИ-1. Подготовить релиз.} Предусловие: администратор зарегистрировал продукт,
его компоненты и среды. Система принимает сведения о сборках. Релиз-менеджер создаёт
черновик, включает версии компонентов, проверяет перечень изменений и фиксирует состав.
Результат: релиз «Сформирован». Неполные сведения о сборке, неизвестный компонент,
пустой состав или две версии одного компонента приводят к отказу в соответствующей операции.
Требования: группы 01 и 02.

\textbf{ВИ-2. Выпустить релиз.} Предусловие: релиз сформирован, задан маршрут с
промышленной средой на последнем этапе. Релиз-менеджер запускает продвижение.
Перед этапом система получает согласования, ожидает освобождения среды и технологического
окна. Затем выполняет шаги по порядку и проверяет КТК. После последнего пройденного этапа
релиз становится выпущенным. При отказе согласующего, ошибке, истечении ожидания или
непройденной КТК продвижение требует вмешательства. Требования: группы 03--07.

\textbf{ВИ-3. Повторить шаг после ошибки.} Предусловие: релиз требует вмешательства,
результат предыдущей попытки установлен, внешнее задание больше не выполняется.
Релиз-менеджер повторяет неуспешный шаг; система сохраняет результаты предшествующих шагов.
После успеха выполнение продолжается с оставшихся шагов, затем КТК проверяется заново.
При неизвестном результате повтор отклоняется. После отклонённого согласования повтор
шага недоступен; новый запрос согласования остаётся за границами версии 1.
Требования: \reqref{SYS_04_09}, \reqref{SYS_04_10}.

\textbf{ВИ-4. Отменить релиз.} Релиз-менеджер отменяет релиз до выпуска.
При отсутствии выполняющегося шага релиз становится отменённым.
Во время выполнения шага система отклоняет отмену; пользователь повторяет её после
получения результата. Новые шаги после отмены не начинаются. Выполненные установки
остаются в учёте. Основание: \reqref{SYS_04_12}.

\textbf{ВИ-5. Откатить выпуск.} Предусловие: в промышленной среде установлен последний
выпущенный релиз. Предыдущий успешный выпуск этого продукта известен.
Релиз-менеджер запускает откат. Система ожидает освобождения среды и разрешённого
времени начала, затем передаёт задание внешней системе развёртывания.
При успехе учёт установок отражает предыдущий состав, откатываемый релиз становится
«Откачен». При ошибке или неизвестном результате операция отката требует вмешательства;
новые работы в этой среде блокируются до установления результата. Требования: группа 08.

\textbf{ВИ-6. Проверить действия.} Предусловие: пользователь аутентифицирован.
Пользователю доступны составы и установки. Аудитор также видит историю изменений:
действия, время и инициаторов. Операции вне прав роли отклоняются.
Связанные требования: \reqref{SYS_01_07} и \reqref{SYS_08_01}; группы 09 и 10.

\textbf{ВИ-7. Продолжить после перезапуска.} Предусловие: продвижение началось,
часть шагов завершена. После перезапуска система восстанавливает подтверждённые
результаты и запрашивает состояние незавершённых внешних заданий. Пока результат
неизвестен, новое задание вместо прежнего не запускается. Требования:
\reqref{SYS_04_10}, \reqref{SYS_11_03}.

Состояния и разрешённые переходы сведены в таблицу~\ref{tab:states}.
При ожидании согласования, свободной среды или окна релиз остаётся «На продвижении».
Откат учитывается как отдельная операция; неудачная попытка не превращает выпуск
в успешно откаченный релиз. Рабочие решения по спорным переходам отмечены в
разделе~\ref{sec:questions}.

\begin{tabl}{|P{3.0cm}|P{3.0cm}|L|}{Переходы состояний релиза}[tab:states]{\textbf{Из состояния} & \textbf{В состояние} & \textbf{Основание}}
Релиза нет & Черновик & Создание: \reqref{SYS_02_01}. \\ \hline
Черновик & Сформирован & Фиксация состава: \reqref{SYS_02_04}. \\ \hline
Сформирован & На продвижении & Запуск: \reqref{SYS_04_01}. \\ \hline
На продвижении & Требует вмешательства & Ошибка, отказ или неизвестный результат: \reqref{SYS_04_08}. \\ \hline
Требует вмешательства & На продвижении & Разрешённый повтор шага: \reqref{SYS_04_09}. \\ \hline
На продвижении & Выпущен & Последний этап пройден: \reqref{SYS_04_11}. \\ \hline
Черновик; Сформирован; На продвижении; Требует вмешательства & Отменён & Отмена при отсутствии выполняющегося шага: \reqref{SYS_04_12}. \\ \hline
Выпущен & Откачен & Успешный откат: \reqref{SYS_08_04}. \\ \hline
\end{tabl}

\section{Архитектура системы и программно-аппаратные интерфейсы}\label{sec:architecture}
\subsection{Границы и функциональный состав}

На системном уровне выделены три части: средства взаимодействия с пользователями,
управление релизами и сохранение данных. Первая часть принимает действия и предоставляет
результаты. Вторая проверяет права и условия продвижения, выдаёт задания внешним системам.
Третья сохраняет каталог, составы, результаты, установки и аудит.
Это функциональное разделение; оно не задаёт число программ, процессов или серверов.

Словесная схема взаимодействия: пользователи обмениваются действиями и результатами
с «Диспетчером релизов»; он получает сведения о сборках, обменивается заданиями и
результатами с системой развёртывания, запрашивает сведения у системы учёта задач
и подтверждение личности у службы аутентификации. Хранилище дистрибутивов обслуживает
внешние системы; «Диспетчер релизов» учитывает ссылки на дистрибутивы.

\subsection{Программные интерфейсы}

В таблице~\ref{tab:interfaces} указаны направления относительно «Диспетчера релизов».
Внешние системы названы по функции. Форматы сообщений и протоколы будут определены
при разработке требований к ПО.

\begin{tabl}{|P{3.1cm}|L|L|}{Обмен с внешними системами}[tab:interfaces]{\textbf{Система} & \textbf{Входящие сведения} & \textbf{Исходящие сведения и событие}}
Система сборки & Компонент, номер версии, ссылка на дистрибутив, контрольная сумма, ревизия исходного кода, время сборки. & Приём сведений после сборки; ответ о регистрации или отказе. \reqref{SYS_01_04}--\reqref{SYS_01_06}. \\ \hline
Система развёртывания & Состояние задания, успех или ошибка; установленные версии; числа выполненных и успешных проверок. & Задание на автоматический шаг либо откат с указанием релиза, среды и состава; запрос состояния незавершённого задания. \reqref{SYS_04_04}, \reqref{SYS_04_05}, \reqref{SYS_04_10}, \reqref{SYS_08_03}. \\ \hline
Система учёта задач & Задачи и дефекты, связанные с версиями; число открытых критических дефектов. & Запрос перечня изменений или данных для КТК. \reqref{SYS_02_07}, \reqref{SYS_06_03}. \\ \hline
Служба аутентификации & Подтверждение личности пользователя либо отказ. & Запрос проверки личности при входе. \reqref{SYS_09_01}. \\ \hline
Служба доставки сообщений & В версии 1 обмен не требуется. & Резерв: уведомление о назначенном действии. Решение о включении --- В-4. \\ \hline
\end{tabl}

Сведения внешних систем учитываются по связанным с ними продукту, компоненту,
релизу, этапу и заданию. При отсутствии достоверного результата система показывает
неопределённость и блокирует зависимое продвижение по требованиям раздела~\ref{sec:functional}.

\subsection{Аппаратное окружение и допущения}

Рабочие места пользователей и вычислительная среда системы связаны сетью.
Специальные аппаратные устройства, датчики и управляющие сигналы проектом не предусмотрены.
Обмен с внешними системами выполняется по сети. В требованиях к ПО будут описаны
защита соединений и проверка подлинности внешних отправителей.
Конкретные протоколы в системных требованиях не заданы.

Для системных испытаний предлагаются имитаторы внешних систем с управляемыми
успешными ответами, ошибками, задержками и потерями связи. Это условие проверки,
а не функция продукта. Восстановление в версии 1 рассматривается при сохранном
хранилище данных после перезапуска; аварийное восстановление разрушенного хранилища
не входит в принятый объём.

\clarify{В-1, преподавателю: достаточны ли функциональное описание и таблица обмена
в разделе 3 как описание архитектуры и программно-аппаратных интерфейсов?
Глубина описания в материалах курса не задана; вопрос снимается после согласования на паре.}

\section{Функциональные требования}\label{sec:functional}

ID имеет вид SYS_ФФ_НН, где ФФ --- номер функции, НН --- номер требования внутри неё.
После ID указана версия требования. Все требования ниже имеют версию 1.
Отклонение операции означает отказ с указанием причины и сохранением данных,
которые эта операция пыталась изменить; запись об отказе в журнале допускается.
Для пользовательских операций действуют ограничения ролей из таблицы~\ref{tab:roles}.

В источниках «черновик SYS-NNN», «резерв R-NN» и «Б.N» ссылаются на часть Б
рабочих материалов команды из списка источников. Черновые номера служат только
для прослеживания переноса; постоянные ID заданы в этом документе.
Допустимый согласующий --- адресат запроса с ролью «Согласующий»,
который не запускал продвижение данного релиза.

\subsection{01. Каталог и версии}

\begin{req}{SYS_01_01}{1}
Система должна предоставлять администратору возможность регистрации программного продукта в каталоге.
\reqsource{П-1; ВИ-1; черновик SYS-001.}
\reqverify{Испытание: зарегистрированный продукт присутствует в каталоге под своим идентификатором.}
\end{req}

\begin{req}{SYS_01_02}{1}
Система должна предоставлять администратору возможность регистрации компонента с привязкой к зарегистрированному программному продукту.
\reqsource{П-1; ВИ-1; черновик SYS-002.}
\reqverify{Испытание: компонент включён в указанный продукт; при неизвестном продукте регистрация отклонена.}
\end{req}

\begin{req}{SYS_01_03}{1}
Система должна предоставлять администратору возможность регистрации среды развёртывания с уникальным идентификатором.
\reqsource{П-2; ВИ-1; черновик SYS-003.}
\reqverify{Испытание: среда доступна для назначения этапу; повторная регистрация того же идентификатора отклонена.}
\end{req}

\begin{req}{SYS_01_04}{1}
Система должна регистрировать версию известного компонента по полученным сведениям о сборке, включающим номер версии, ссылку на дистрибутив, контрольную сумму, ревизию исходного кода и время сборки.
\reqsource{П-1, П-9; ВИ-1; черновик SYS-004.}
\reqverify{Испытание: передать полный набор сведений; зарегистрированная версия содержит значения из сообщения.}
\reqnote{Время регистрации задано отдельно в \reqref{SYS_11_01}.}
\end{req}

\begin{req}{SYS_01_05}{1}
Система должна отклонять сведения о сборке неизвестного компонента или сведения без хотя бы одного обязательного значения, перечисленного в \reqref{SYS_01_04}.
\reqsource{П-1, П-9; ВИ-1; резерв R-04 и состав данных Б.8.}
\reqverify{Испытание: передать неизвестный компонент и поочерёдно исключить каждое обязательное значение; во всех случаях новая версия отсутствует, сообщена причина отказа.}
\end{req}

\begin{req}{SYS_01_06}{1}
Система должна сохранять единственную неизменную запись версии для каждой пары «компонент, номер версии»: повтор идентичных сведений подтверждает существующую запись, а отличающиеся сведения отклоняются.
\reqsource{П-1; ВИ-1; резерв R-03.}
\reqverify{Испытание: два одинаковых сообщения дают одну запись; сообщение с тем же номером и другой контрольной суммой отклонено, прежняя запись сохранена.}
\end{req}

\begin{req}{SYS_01_07}{1}
Система должна предоставлять пользователю с любой из семи ролей просмотр каталога продуктов, компонентов, версий и сред, а также состава и текущего состояния релизов.
\reqsource{П-1, П-2, П-8; ВИ-6; матрица прав Б.7.}
\reqverify{Испытание: каждой из семи ролей доступен просмотр сохранённых сведений. Пользователь без роли получает отказ.}
\end{req}

\subsection{02. Релизы}

\begin{req}{SYS_02_01}{1}
Система должна предоставлять релиз-менеджеру возможность создания релиза зарегистрированного продукта в состоянии «Черновик».
\reqsource{П-1; ВИ-1; черновик SYS-005.}
\reqverify{Испытание: созданный релиз связан с выбранным продуктом и имеет состояние «Черновик».}
\end{req}

\begin{req}{SYS_02_02}{1}
Система должна предоставлять релиз-менеджеру возможность включения зарегистрированной версии компонента в состав релиза в состоянии «Черновик».
\reqsource{П-1; ВИ-1; черновик SYS-006.}
\reqverify{Испытание: выбранная допустимая версия появилась в составе черновика.}
\end{req}

\begin{req}{SYS_02_03}{1}
Система должна предоставлять релиз-менеджеру возможность удаления версии компонента из состава релиза в состоянии «Черновик».
\reqsource{П-1; ВИ-1; возможность изменения черновика в Б.6.}
\reqverify{Испытание: удалённая версия отсутствует в составе, остальные версии сохранены.}
\end{req}

\begin{req}{SYS_02_04}{1}
Система должна по запросу релиз-менеджера фиксировать непустой состав черновика с переводом релиза в состояние «Сформирован».
\reqsource{П-1; ВИ-1; черновик SYS-007.}
\reqverify{Испытание: непустой допустимый состав зафиксирован; пустой черновик остаётся черновиком с отказом в фиксации.}
\end{req}

\begin{req}{SYS_02_05}{1}
Система должна отклонять любое изменение состава релиза после его фиксации.
\reqsource{П-1; ВИ-1; черновик SYS-008.}
\reqverify{Испытание: попытки включить, заменить и удалить версии после фиксации отклонены; состав прежний.}
\end{req}

\begin{req}{SYS_02_06}{1}
Система должна отклонять включение в состав релиза неизвестной версии, версии компонента другого продукта или второй версии уже включённого компонента.
\reqsource{П-1; ВИ-1; модель Б.5 и резерв R-02.}
\reqverify{Испытание: каждый из трёх недопустимых вариантов отклонён; состав черновика сохранён.}
\end{req}

\begin{req}{SYS_02_07}{1}
Система должна формировать перечень изменений релиза из задач и исправленных дефектов, связанных во внешней системе учёта задач с версиями его текущего состава.
\reqsource{П-1, П-9; ВИ-1; черновик SYS-009.}
\reqverify{Испытание: задать связи задач с версиями; перечень содержит ровно связанные с составом задачи и исправленные дефекты, без повторов.}
\end{req}

\begin{req}{SYS_02_08}{1}
Система должна обозначать перечень изменений как неполный, если сведения внешней системы учёта задач для хотя бы одной версии состава недоступны.
\reqsource{П-1, П-9; ВИ-1; резерв R-22.}
\reqverify{Испытание: прекратить ответы по одной версии; перечень содержит явную отметку о неполноте, а не обозначение пустого полного списка.}
\end{req}

\subsection{03. Маршрут продвижения}

\begin{req}{SYS_03_01}{1}
Система должна предоставлять релиз-менеджеру возможность задания упорядоченного списка этапов маршрута продукта с указанием зарегистрированной среды каждого этапа.
\reqsource{П-2; ВИ-2; черновик SYS-010.}
\reqverify{Испытание: задать три этапа; сохранены их порядок и среды; неизвестная среда отклонена.}
\end{req}

\begin{req}{SYS_03_02}{1}
Система должна предоставлять релиз-менеджеру возможность включения в этап автоматического шага установки состава релиза либо автоматической проверки с указанием положительного времени ожидания результата в секундах.
\reqsource{П-2, П-9; ВИ-2; черновик SYS-011, SYS-017.}
\reqverify{Испытание: шаг с действием «установка» и ожиданием 60 с принят; нулевое и отрицательное ожидание отклонены.}
\end{req}

\begin{req}{SYS_03_03}{1}
Система должна предоставлять релиз-менеджеру возможность включения в этап ручного шага с назначением исполнителя: инженера по тестированию для непромышленной среды либо инженера эксплуатации для промышленной.
\reqsource{П-2, П-8; ВИ-2; черновик SYS-012; матрица Б.7.}
\reqverify{Испытание: шаг с подходящим исполнителем принят; назначение пользователя без требуемой роли отклонено.}
\end{req}

\begin{req}{SYS_03_04}{1}
Система должна предоставлять релиз-менеджеру возможность задания порядка шагов внутри этапа.
\reqsource{П-2; ВИ-2; порядок выполнения из Б.6.}
\reqverify{Испытание: сохранить три шага в заданном порядке; при чтении маршрута порядок совпадает с заданным.}
\end{req}

\begin{req}{SYS_03_05}{1}
Система должна использовать для каждого запущенного продвижения маршрут, состав шагов, КТК и правила согласования в том виде, в каком они были на момент запуска.
\reqsource{П-2, П-3; ВИ-2; резерв R-06.}
\reqverify{Испытание: изменить маршрут продукта после запуска; действующий релиз проходит прежний маршрут, следующий использует новый.}
\end{req}

\subsection{04. Продвижение релиза}

\begin{req}{SYS_04_01}{1}
Система должна по запросу релиз-менеджера запускать продвижение сформированного релиза с переводом в состояние «На продвижении» при выполнении всех условий:
\begin{itemize}
\item маршрут содержит хотя бы один этап, каждый этап содержит шаги;
\item последний этап относится к промышленной среде и включает установку состава релиза;
\item для каждого правила согласования число допустимых согласующих не меньше N.
\end{itemize}
\reqsource{П-2, П-3; ВИ-2; черновик SYS-013.}
\reqverify{Испытание: допустимый релиз начинает ожидание допуска к первому этапу; запуск черновика, пустого маршрута, этапа без шагов или правила без достаточного числа согласующих отклонён.}
\end{req}

\begin{req}{SYS_04_02}{1}
Система должна начинать выполнение каждого этапа после первого только после перехода предыдущего этапа в состояние «Пройден».
\reqsource{П-2; ВИ-2; черновик SYS-014.}
\reqverify{Испытание: оставить предыдущий этап незавершённым или завершить с ошибкой; следующий этап не начат.}
\end{req}

\begin{req}{SYS_04_03}{1}
Система должна выполнять шаги этапа по порядку, начиная следующий шаг только после успеха предыдущего.
\reqsource{П-2; ВИ-2; последовательность шагов Б.6.}
\reqverify{Испытание: при незавершённом либо ошибочном предыдущем шаге следующий не начинается; после успеха начинается следующий по порядку.}
\end{req}

\begin{req}{SYS_04_04}{1}
Система должна при начале автоматического шага передавать внешней системе развёртывания задание с указанием релиза, среды, действия и зафиксированного состава релиза.
\reqsource{П-9; ВИ-2; черновик SYS-015.}
\reqverify{Испытание: имитатор получает задание, значения которого совпадают с текущим этапом и составом релиза.}
\end{req}

\begin{req}{SYS_04_05}{1}
Система должна отражать успешный или ошибочный результат автоматического шага по результату соответствующего внешнего задания.
\reqsource{П-2, П-9; ВИ-2; черновик SYS-015.}
\reqverify{Испытание: ответы об успехе и ошибке изменяют состояние своего шага; результат другого задания этот шаг не изменяет.}
\end{req}

\begin{req}{SYS_04_06}{1}
Система должна принимать результат ручного шага «выполнен» или «не выполнен» только от назначенного исполнителя во время выполнения этого шага.
\reqsource{П-2, П-8; ВИ-2; черновик SYS-016.}
\reqverify{Испытание: допустимый результат исполнителя отображается как «Успешно» либо «Ошибка»; результат другого пользователя или для ожидающего шага отклонён.}
\end{req}

\begin{req}{SYS_04_07}{1}
Система должна переводить этап в состояние «Пройден» после успешного выполнения всех его шагов и выполнения всех критериев назначенной КТК; для этапа без КТК достаточно успешного выполнения всех шагов.
\reqsource{П-2, П-3; ВИ-2; черновик SYS-025 и порядок Б.6.}
\reqverify{Испытание: этап пройден, когда все шаги успешны и КТК выполнена. Ошибка шага или невыполненный критерий препятствуют прохождению. Этап без КТК завершается после последнего успешного шага.}
\end{req}

\begin{req}{SYS_04_08}{1}
Система должна переводить продвигаемый релиз в состояние «Требует вмешательства» при ошибке шага, невыполненной КТК, отклонённом согласовании либо истечении заданного времени ожидания результата автоматического шага.
\reqsource{П-2, П-3; ВИ-2; черновик SYS-017.}
\reqverify{Испытание: воспроизвести каждое из четырёх условий; релиз требует вмешательства, следующий шаг не начинается.}
\reqnote{Время ожидания отсчитывается от передачи задания. Отсутствие данных для КТК трактуется по \reqref{SYS_06_04}.}
\end{req}

\begin{req}{SYS_04_09}{1}
Система должна по запросу релиз-менеджера повторять выбранный шаг непройденного
текущего этапа релиза в состоянии «Требует вмешательства» с возвратом в состояние
«На продвижении» при выполнении всех условий:
\begin{itemize}
\item результат прежнего внешнего задания установлен;
\item среда доступна для этого релиза;
\item согласование этапа сохраняет силу либо не требуется.
\end{itemize}
Повтор выполняется с прежними параметрами; результаты предшествующих шагов сохраняются,
результаты выбранного и следующих шагов остаются только в истории. После повторного
прохождения оставшихся шагов КТК проверяется заново.
\reqsource{П-2, П-3; ВИ-3; черновик SYS-018.}
\reqverify{Испытание: повторить ошибочный шаг — предшествующие шаги сохранены, выполнение продолжается от выбранного; после повторения КТК проверяется заново; при отклонённом согласовании повтор недоступен.}
\end{req}

\begin{req}{SYS_04_10}{1}
Система должна при неизвестном результате внешнего задания запрашивать его состояние без выдачи повторного задания на то же действие до установления результата.
\reqsource{П-2, П-4, П-9; ВИ-3, ВИ-7; уточнение черновика SYS-017, SYS-040.}
\reqverify{Испытание: потерять ответ, затем перезапустить систему; имитатор получает запрос состояния, число заданий на исходное действие остаётся равным одному до получения результата.}
\end{req}

\begin{req}{SYS_04_11}{1}
Система должна переводить релиз в состояние «Выпущен» после перехода последнего этапа его маршрута в состояние «Пройден».
\reqsource{П-2; ВИ-2; черновик SYS-019.}
\reqverify{Испытание: пройти весь маршрут — релиз выпущен; до прохождения последнего этапа состояние «Выпущен» отсутствует.}
\end{req}

\begin{req}{SYS_04_12}{1}
Система должна по запросу релиз-менеджера отменять релиз из одного из состояний:
«Черновик», «Сформирован», «На продвижении», «Требует вмешательства».
Отмена допускается только при отсутствии выполняющегося шага и внешнего задания
с неизвестным результатом; итоговое состояние --- «Отменён».
\reqsource{П-2, П-4; ВИ-4; черновик SYS-020 и вопрос Б.12.10.}
\reqverify{Испытание: отменить релиз в каждом из допустимых состояний — дальнейшие шаги не запускаются; во время внешнего или ручного шага отказ, продвижение сохранено.}
\reqnote{Рабочее правило версии 1; подтверждение требуется по В-3.}
\end{req}

\begin{req}{SYS_04_13}{1}
Система должна отклонять операции, вызывающие переход состояния релиза, отсутствующий в таблице~\ref{tab:states}.
\reqsource{П-2; ВИ-1--ВИ-5; модель Б.6.}
\reqverify{Испытание: попытаться вернуть сформированный релиз в черновик, запустить отменённый и отменить выпущенный; каждая операция отклонена.}
\end{req}

\subsection{05. Доступ к средам}

\begin{req}{SYS_05_01}{1}
Система должна предоставлять среду одновременно только одному этапу или операции отката,
занимая её при начале работы и освобождая после завершения работы с установленными
результатами внешних заданий и подтверждённым состоянием среды.
При неизвестном результате или неподтверждённом состоянии занятость сохраняется;
ожидание согласования или окна само по себе среду не занимает.
\reqsource{П-4; ВИ-2, ВИ-5; черновик SYS-021.}
\reqverify{Испытание: одновременно допустить два релиза либо релиз и откат в одну среду; выполняется одна работа, вторая ожидает, в том числе при потере ответа о первой.}
\end{req}

\begin{req}{SYS_05_02}{1}
Система должна предоставлять инженеру эксплуатации возможность задания технологических окон среды как интервалов с датой, временем начала, временем окончания и часовым смещением, в которых начало предшествует окончанию.
\reqsource{П-4; ВИ-2; черновик SYS-022.}
\reqverify{Испытание: сохранить интервал, в том числе через полночь; равные границы и окончание раньше начала отклонены.}
\end{req}

\begin{req}{SYS_05_03}{1}
Система должна начинать этап, его повтор или откат в среде с заданными технологическими окнами только при времени начала не раньше начала и строго раньше окончания хотя бы одного окна.
При отсутствии окон ограничение по времени отсутствует. Завершение окна не прерывает
уже начатую работу.
\reqsource{П-4; ВИ-2, ВИ-5; черновик SYS-023.}
\reqverify{Испытание: проверить момент до начала, на начале и на окончании окна; старт допустим только на начале. При отсутствии окон ограничение по времени отсутствует.}
\reqnote{Рабочее правило по В-2 и В-3.}
\end{req}

\subsection{06. Контрольные точки качества}

\begin{req}{SYS_06_01}{1}
Система должна предоставлять релиз-менеджеру возможность назначения этапу одной КТК с непустым набором критериев: минимальной долей успешных автоматических проверок и/или максимальным числом открытых критических дефектов.
\reqsource{П-3; ВИ-2; черновик SYS-024, SYS-025.}
\reqverify{Испытание: назначить каждый тип критерия и оба вместе; сохранённый набор относится к выбранному этапу.}
\end{req}

\begin{req}{SYS_06_02}{1}
Система должна отклонять критерий КТК с порогом доли вне диапазона от 0 до 100\% включительно или с порогом числа дефектов, не являющимся целым неотрицательным числом.
\reqsource{П-3; ВИ-2; проверяемость критериев Б.9.}
\reqverify{Испытание: пороги 0 и 100\%, а также 0 дефектов принимаются; −1\%, 101\%, −1 и 0,5 дефекта отклоняются.}
\end{req}

\begin{req}{SYS_06_03}{1}
Система должна определять выполнение КТК после завершения шагов текущей попытки этапа: доля успешных проверок по его отчётам должна быть не ниже заданного порога, а число открытых критических дефектов релиза на момент проверки — не выше заданного порога.
Учитываются только назначенные критерии. Доля равна числу успешных проверок,
делённому на общее число, умноженному на 100; сравнение выполняется без округления.
\reqsource{П-3, П-9; ВИ-2, ВИ-3; черновик SYS-025.}
\reqverify{Испытание: при пороге 95\% результаты 95 из 100 проходят, 94 из 100 не проходят; при пороге 0 дефектов число 0 проходит, 1 не проходит; при двух критериях нужны оба.}
\end{req}

\begin{req}{SYS_06_04}{1}
Система должна считать КТК невыполненной при отсутствии данных хотя бы для одного назначенного критерия, нулевом общем числе проверок для критерия доли либо противоречивом отчёте о проверках.
Противоречивым считается отчёт с отрицательными или нецелыми счётчиками либо
числом успешных проверок больше общего.
\reqsource{П-3, П-9; ВИ-2; ненормальные условия для SYS-025.}
\reqverify{Испытание: недоступный учёт задач, отсутствующий отчёт, 0 проверок и число успешных больше общего каждый раз препятствуют прохождению КТК.}
\end{req}

\subsection{07. Согласования}

\begin{req}{SYS_07_01}{1}
Система должна предоставлять релиз-менеджеру возможность назначения этапу правила «N из M», где M — число различных пользователей с ролью «Согласующий», а N — целое число от 1 до M включительно.
\reqsource{П-3; ВИ-2; черновик SYS-026.}
\reqverify{Испытание: правило «2 из 3» сохраняется; пустой список, N=0, N>M и повтор одного пользователя в списке отклоняются.}
\end{req}

\begin{req}{SYS_07_02}{1}
Система должна допускать релиз к этапу с правилом согласования после получения N положительных решений разных допустимых согласующих при отсутствии отклоняющего решения в открытом запросе.
При достижении N запрос закрывается с результатом «согласовано».
Для начала этапа также действуют требования к предыдущему этапу и доступу к среде.
\reqsource{П-3; ВИ-2; черновик SYS-027.}
\reqverify{Испытание: для «2 из 3» одного голоса недостаточно, двух достаточно; повтор голоса первого пользователя не увеличивает число голосов.}
\end{req}

\begin{req}{SYS_07_03}{1}
Система должна принимать одно решение «согласовано» или «отклонено» от каждого адресата открытого запроса на согласование.
\reqsource{П-3; ВИ-2; черновик SYS-028.}
\reqverify{Испытание: адресат с нужной ролью принимает одно решение; повтор, решение постороннего или решение по закрытому запросу отклонены.}
\end{req}

\begin{req}{SYS_07_04}{1}
Система должна отклонять решение по согласованию релиза от пользователя, запустившего его продвижение, независимо от совмещения ролей.
\reqsource{П-3, П-8; ВИ-2; черновик SYS-029.}
\reqverify{Испытание: инициатор с ролями релиз-менеджера и согласующего пытается согласовать или отклонить свой релиз; оба решения отклонены.}
\end{req}

\begin{req}{SYS_07_05}{1}
Система должна закрывать открытый запрос на согласование с результатом «отклонено» после первого допустимого отклоняющего решения.
\reqsource{П-3; ВИ-2; правило вето Б.6.}
\reqverify{Испытание: при «2 из 3» после одного положительного и одного отрицательного решения запрос отклонён, релиз требует вмешательства, последующее положительное решение не принимается.}
\reqnote{Рабочее правило вето до закрытия запроса; требует решения команды по В-2.}
\end{req}

\subsection{08. Учёт установок и откат}

\begin{req}{SYS_08_01}{1}
Система должна предоставлять пользователю с любой из семи ролей сведения о подтверждённых установленных версиях компонентов каждой среды с указанием времени и основания последнего подтверждения.
\reqsource{П-6; ВИ-5, ВИ-6; черновик SYS-030.}
\reqverify{Испытание: после подтверждённой установки просмотр среды показывает версии из результата задания и время подтверждения; при отсутствии подтверждений это явно указано.}
\end{req}

\begin{req}{SYS_08_02}{1}
Система должна обновлять сведения об установленных версиях только по подтверждённым результатам установки или отката, сохраняя пометку «текущее состояние не подтверждено» для среды с ошибочным либо неизвестным результатом изменения её состава.
\reqsource{П-5, П-6; ВИ-2, ВИ-5; черновик SYS-030, SYS-032.}
\reqverify{Испытание: успешная установка или откат меняет учёт; ошибка и потеря ответа оставляют последние подтверждённые сведения с пометкой, а не выдают целевой состав за установленный.}
\end{req}

\begin{req}{SYS_08_03}{1}
Система должна по запросу релиз-менеджера передавать задание на откат последнего выпущенного релиза к предшествующему успешному выпуску того же продукта, если текущая установка соответствует откатываемому релизу и выполнены условия доступа к промышленной среде.
\reqsource{П-5, П-4; ВИ-5; черновик SYS-031.}
\reqverify{Испытание: после выпусков R1 и R2 задание на откат R2 содержит состав R1; без предшествующего выпуска или при несовпадении текущего состава запрос отклонён, при занятой среде или закрытом окне ожидает.}
\reqnote{См. \reqref{SYS_08_02}, \reqref{SYS_05_01} и \reqref{SYS_05_03}.}
\end{req}

\begin{req}{SYS_08_04}{1}
Система должна переводить откатываемый релиз в состояние «Откачен» после подтверждения внешней системой успешного восстановления всего целевого состава.
\reqsource{П-5; ВИ-5; черновик SYS-032.}
\reqverify{Испытание: подтверждён весь состав R1 — R2 откачен; частичное восстановление или отсутствие результата не дают этого состояния.}
\end{req}

\begin{req}{SYS_08_05}{1}
Система должна обозначать операцию отката как требующую вмешательства при ошибке или истечении заданного времени ожидания её результата, сохраняя состояние релиза «Выпущен».
Время ожидания задаёт релиз-менеджер при запросе как положительное число секунд;
запрос с нулевым или отрицательным временем отклоняется.
\reqsource{П-5, П-6; ВИ-5; вопрос Б.12.7.}
\reqverify{Испытание: вернуть ошибку или потерять ответ до истечения ожидания; операция требует вмешательства, релиз не помечен откаченным, состояние среды показано как неподтверждённое.}
\reqnote{Восстановление после частичного отката требует решения В-3. До подтверждения состава среда остаётся занятой по \reqref{SYS_05_01}.}
\end{req}

\subsection{09. Управление доступом}

\begin{req}{SYS_09_01}{1}
Система должна выполнять пользовательские операции с данными только при подтверждённой внешней службой аутентификации личности пользователя.
\reqsource{П-8; ВИ-1--ВИ-6; черновик SYS-033.}
\reqverify{Испытание: подтверждённая личность допускается к проверке прав; отказ и недоступность службы препятствуют входу и выполнению операций.}
\reqnote{Вход предшествует операциям с данными. Внешняя аутентификация предложена в В-4.}
\end{req}

\begin{req}{SYS_09_02}{1}
Система должна предоставлять администратору возможность назначения пользователю одной или нескольких ролей из таблицы~\ref{tab:roles}.
\reqsource{П-8; ВИ-6; черновик SYS-034.}
\reqverify{Испытание: назначить одну роль, затем вторую; сохранены обе роли; неизвестная роль отклонена.}
\end{req}

\begin{req}{SYS_09_03}{1}
Система должна отклонять операцию пользователя, не разрешённую ни одной из его ролей по таблице~\ref{tab:roles}, с учётом ограничений назначенного исполнителя и инициатора продвижения.
\reqsource{П-8; ВИ-6; черновик SYS-035.}
\reqverify{Испытание: проверить разрешения каждой роли по таблице; аудитор не запускает продвижение, администратор без роли аудитора не читает журнал, совмещение ролей сохраняет запрет самосогласования.}
\end{req}

\subsection{10. Журнал аудита}

\begin{req}{SYS_10_01}{1}
Система должна регистрировать каждое изменение своих предметных данных в журнале аудита с указанием даты и времени, инициатора, объекта, действия и результата изменения.
\reqsource{П-7; ВИ-6; черновик SYS-036.}
\reqverify{Испытание: изменить каталог, роли, состав, маршрут, состояние релиза, решение и установку; каждому изменению соответствует запись со всеми значениями.}
\reqnote{Для автоматического изменения инициатор — система, с указанием вызвавшего его события. Добавление записи аудита не порождает запись о самой себе.}
\end{req}

\begin{req}{SYS_10_02}{1}
Система должна отклонять редактирование и удаление существующих записей журнала аудита через любые предоставляемые ею операции независимо от роли пользователя.
\reqsource{П-7; ВИ-6; черновик SYS-037.}
\reqverify{Испытание: попытки изменить и удалить запись с каждой ролью, включая администратора, отклонены; исходная запись сохранена.}
\reqnote{Граница требования — операции системы; защита от физического разрушения хранилища в версию 1 не входит.}
\end{req}

\begin{req}{SYS_10_03}{1}
Система должна предоставлять аудитору возможность просмотра журнала аудита в хронологическом порядке.
\reqsource{П-7; ВИ-6; черновик SYS-038.}
\reqverify{Испытание: аудитор видит зарегистрированные изменения с их временем и инициаторами в порядке времени.}
\end{req}

\section{Нефункциональные требования}\label{sec:nonfunctional}

Показатели этого раздела являются предложением версии 1. Условия нагрузочной проверки подлежат согласованию по В-5. Требования к содержимому и неизменяемости журнала заданы в группе 10.

\begin{req}{SYS_11_01}{1}
Система должна завершать регистрацию версии компонента не позднее чем через 60 с после получения полных допустимых сведений о сборке при нагрузке до 20 одновременно продвигаемых релизов.
\reqsource{П-1, П-9; ВИ-1; черновик SYS-004.}
\reqverify{Испытание: при указанной нагрузке измерить интервал от приёма сведений до доступности версии в каталоге; интервал не превышает 60 с.}
\reqnote{Значения предложены в черновике и требуют подтверждения по В-5.}
\end{req}

\begin{req}{SYS_11_02}{1}
Система должна поддерживать одновременное продвижение 20 релизов при независимых целевых средах и доступных внешних системах.
\reqsource{П-2; ВИ-2; черновик SYS-039.}
\reqverify{Испытание: запустить 20 релизов разных продуктов с независимыми средами; все проходят маршрут, ни один не отклонён из-за числа релизов.}
\reqnote{Конкуренция за одну среду проверяется отдельно по \reqref{SYS_05_01}; 20 релизов — проверяемая нагрузка, а не запрет на большее число.}
\end{req}

\begin{req}{SYS_11_03}{1}
Система должна после перезапуска при сохранном хранилище восстанавливать все подтверждённые до перезапуска состояния релизов, результаты шагов, решения согласующих и записи аудита до возобновления продвижения.
\reqsource{П-2, П-7; ВИ-7; черновик SYS-040.}
\reqverify{Испытание: перезапустить систему на автоматическом шаге, при ожидании согласования и после фиксации результата; потерь подтверждённых записей — 0, повторных заданий до выяснения результата — 0.}
\reqnote{Незавершённые задания проверяются по \reqref{SYS_04_10}. Время запуска зависит от согласуемой испытательной среды, В-5.}
\end{req}

\begin{req}{SYS_11_04}{1}
Система должна предоставлять результат просмотра каталога, состава и состояния релиза либо журнала аудита не позднее чем через 3 с после получения запроса при одновременной работе 30 пользователей и продвижении 20 релизов.
\reqsource{П-1, П-2, П-7; ВИ-6; резерв R-20.}
\reqverify{Испытание: в течение согласованного нагрузочного прогона 30 пользователей выполняют по одному запросу просмотра одновременно с продвижением 20 релизов; каждый ответ укладывается в 3 с.}
\reqnote{Время измеряется на границе системы без сетевой задержки до пользователя. Объём данных и длительность прогона требуют решения В-5; до него критерий неполон.}
\end{req}

\section{Ограничения с обоснованием}
\subsection{Границы версии 1}

Приняты следующие ограничения учебного проекта:
\begin{enumerate}
\item Один маршрут на продукт; этапы и шаги идут последовательно.
Это ограничивает число сценариев конкуренции и позволяет проверить продвижение целиком.
Параллельная работа разных релизов допускается в разных средах.
\item Семь фиксированных ролей. Конструктор прав и разделение одной установки системы
на независимые организации не входят в версию 1.
\item Сборка и выполнение установок переданы внешним системам. Собственные средства
сборки, хранения дистрибутивов и развёртывания не разрабатываются.
\item Откат относится к промышленной среде и предшествующему успешному выпуску
того же продукта. Произвольный выбор старой версии и автоматическое исправление
частичной неудачной установки не включены.
\item Повтор после ошибки использует прежние параметры шага. Новый запрос согласования
после отказа, ручной обход КТК, приостановка продвижения и уведомления внешней службой
остаются резервными функциями.
\item Учёт установок основан на подтверждениях внешней системы. Изменение среды
в обход управляемых заданий не отслеживается автоматически.
\end{enumerate}

Предполагается, что система развёртывания позволяет сопоставить результат с заданием
и получить состояние ранее выданного задания. При отсутствии этой возможности повтор
после потери ответа остаётся заблокированным; это подлежит проверке на имитаторе.

\subsection{Вопросы для согласования}\label{sec:questions}

Вопросы ниже возникли из неполных или расходящихся исходных материалов. Рабочие
варианты позволяют прочитать версию 1 целиком, но до ответа команды соответствующие
требования не считаются согласованными. В-1 приведён в разделе~\ref{sec:architecture}.

\clarify{В-2, команде: подтвердить правило вето до закрытия согласования и трактовку
технологического окна как ограничения начала работ. В исходном черновике вето
отмечено вопросом, а определения окна расходятся. После решения проверить
SYS_07_02, SYS_07_05, SYS_05_03 и глоссарий.}

\clarify{В-3, команде: подтвердить отказ в отмене во время выполняющегося шага,
ожидание технологического окна для отката и сохранение статуса «Выпущен» при
неудаче отката. Определить, как эксплуатация подтверждает состояние среды после
частичного отката. В исходных материалах эти решения отсутствуют; после ответа
уточнить SYS_04_12, SYS_08_03--SYS_08_05 и сценарии ВИ-4, ВИ-5.}

\clarify{В-4, команде: подтвердить внешнюю аутентификацию, использование имитаторов
и резервный статус доставки уведомлений. В README перечислены пять внешних систем,
но уведомления вынесены в резерв, а способ аутентификации оставлен вопросом.
После решения уточнить таблицу обмена и SYS_09_01.}

\clarify{В-5, команде: обосновать показатели 60 с, 20 релизов и 3 с при 30 пользователях;
задать объём каталога и журнала, характеристики испытательной среды, длительность
нагрузочного прогона и допустимое время восстановления. Числа взяты из предложений
черновика и резерва, измерений нет. После согласования дополнить критерии группы 11.}

\clarify{В-6, команде: утвердить рабочее название «Диспетчер релизов» и решить,
являются ли ПСИ отдельным этапом с собственной средой или ручным шагом тестового этапа.
В версии 1 допускаются оба маршрута; вопрос снимается выбором демонстрационного сценария.}

\subsection{Результат самопроверки}

В таблице~\ref{tab:review} указаны свойства, которые остаются под вопросом после
проверки версии 1 по слайдам 22--23. Неделимость, наличие источника и способа проверки
проверены для каждого требования. Полнота и непротиворечивость оцениваются для
документа в целом. Приведённые способы проверки описывают будущие испытания системы;
сама система в рамках подготовки документа не испытывалась.

\begin{tabl}{|P{3.8cm}|P{3.0cm}|L|}{Оставшиеся вопросы по свойствам требований}[tab:review]{\textbf{Требования или раздел} & \textbf{Свойство} & \textbf{Причина и условие закрытия}}
SYS_05_03; SYS_07_02; SYS_07_05 & Корректность & Правила окна и вето требуют решения В-2. \\ \hline
SYS_04_12; SYS_08_03--SYS_08_05 & Корректность; полнота спецификации & Не подтверждены отмена и восстановление после неудачного отката. Закрыть В-3. \\ \hline
Раздел~\ref{sec:architecture}; SYS_09_01 & Полнота спецификации; корректность & Согласовать глубину архитектуры и внешние зависимости по В-1, В-4. \\ \hline
SYS_11_01--SYS_11_04 & Корректность; верифицируемость & Нужны обоснование показателей и воспроизводимые условия нагрузки по В-5. \\ \hline
\end{tabl}

\begin{changelog}
\change{1}{03.10.2026}{Первая версия системных требований; общий глоссарий}{Брылкина К. К.}{Задание к 03.10.2026}
\end{changelog}

\end{document}
